\section{Nonflat geometric base change of the ordered section complex}

Let $R\to S$ be an arbitrary map of commutative rings, let
$X$ be a separated scheme with $p:X\to\operatorname{Spec}R$, and let $M$ be a
quasicoherent module on $X$.  Put
$X_S=X\times_{\operatorname{Spec}R}\operatorname{Spec}S$, with
projections $g:X_S\to X$ and $q:X_S\to\operatorname{Spec}S$.
Fix a finite linearly ordered family of affine opens $U_i$ of $X$.
Write $C_R^\bullet(M,U)$ for its original ordered section complex,
whose degree-$a$ term is the product of the original sections on
intersections indexed by strictly increasing $(a+1)$-tuples.  Its
coefficient action is the literal $R$ action obtained from the
canonical isomorphism $R\simeq\Gamma(\operatorname{Spec}R,\mathcal O)$.
The analogous complex on $X_S$ has the literal original $S$ action.

\begin{lemma}[The original affine section comparison]
\label{lem:kn-geometric-cech-affine-section}
\lean{EnvelopeProperty.SchemeModuleFlatBaseChangeAffineOpen.isIso_comparison,EnvelopeProperty.SchemeModuleFlatBaseChangeMap.comparison_tmul,EnvelopeProperty.SchemeModuleFlatBaseChangeMap.comparison}
\leanok
For every affine open $V\subseteq X$, the canonical map
\[
 S\otimes_R\Gamma(V,M)\longrightarrow
 \Gamma(g^{-1}V,g^*M),\qquad s\otimes m\longmapsto s\,g^*m,
\]
is an isomorphism.  These canonical maps commute with all original
open restrictions, without a flatness assumption on $R\to S$.
\end{lemma}
\begin{proof}
Write $V=\operatorname{Spec}A$ and let $P=\Gamma(V,M)$.
Quasicoherence identifies the restriction of $M$ with the associated
sheaf of $P$.  Its actual inverse image is
$\operatorname{Spec}(S\otimes_R A)$, and the pulled module is the
associated sheaf of $(S\otimes_R A)\otimes_A P$.
Tensor associativity gives the canonical isomorphism with
$S\otimes_R P$, sending every $s\otimes m$ to the original scalar
multiple of the original pullback-unit section.  The pullback
adjunction unit is a natural sheaf morphism.  Its naturality and the
tensor universal property prove restriction compatibility of these
same canonical maps.  No assertion about tensoring kernels is used.
\end{proof}

\begin{theorem}[The whole original ordered complex commutes with base change]
\label{thm:kn-geometric-cech-base-change}
\lean{EnvelopeProperty.Vanishing.GeometricCechBaseChange.complexIso,EnvelopeProperty.Vanishing.GeometricCechBaseChange.termEquiv_tmul}
\leanok
\uses{lem:kn-geometric-cech-affine-section}
There is a canonical isomorphism of actual cochain complexes of
$S$-modules
\[
 S\otimes_R C_R^\bullet(M,U)\ \simeq\
 C_S^\bullet(g^*M,(g^{-1}U_i)_i).
\]
For every ordered tuple, the degreewise map retains precisely the
canonical scalar and pullback-unit formula of the preceding lemma.
Neither $S$ nor $M$ is assumed flat over $R$.
\end{theorem}
\begin{proof}
Separatedness makes every finite intersection of the original affine
opens affine.  In each degree, the ordered index set is finite.
Tensoring with $S$ commutes with its finite product, since a finite
product of modules is also their finite direct sum.  Apply the
canonical affine-section isomorphism to every factor.  Inverse image
preserves these finite intersections, giving exactly the factors of
the actual inverse-image ordered complex.

For each deleted entry in an ordered tuple, the original coface is
restriction to its larger intersection.  The canonical section maps
commute with this restriction by the preceding lemma.  The finite
product tensor isomorphism retains every component on a pure tensor,
so the assembled degreewise isomorphisms commute with each original
coface.  Add the finitely many cofaces with their original signs
$(-1)^j$.  Linearity gives compatibility with the whole alternating
differential.  Thus the degreewise isomorphisms form an isomorphism of
the original cochain complexes.  All coefficient changes use the
canonical $R$ and $S$ global-section isomorphisms, so the comparison
retains the specified original ring actions.
\end{proof}
